\documentclass[sigconf,nonacm]{acmart}

\settopmatter{printacmref=false}

\usepackage{amsmath}
% amssymb omitted: acmart's math setup already defines AMS symbols (\Bbbk etc.); loading amssymb again triggers ``already defined'' errors.
\usepackage{booktabs}
\usepackage{graphicx}
\usepackage{url}
\usepackage{microtype}
\usepackage{array}
\usepackage{tcolorbox}
\tcbuselibrary{listings,breakable,skins}
% Avoid writing colored prompt titles into PDF bookmarks (math, \\, etc.).
\tcbset{bookmark=false}
\definecolor{DefnBoxGray}{gray}{0.94}
% Courier (pcr): avoids inconsolata (zi4) missing italic in listings.
\newcommand{\PromptMono}{\footnotesize\fontencoding{T1}\fontfamily{pcr}\selectfont}
\lstdefinestyle{promptlist}{
  basicstyle=\PromptMono,
  columns=fullflexible,
  breaklines=true,
  breakautoindent=false,
  tabsize=2,
  showstringspaces=false,
  frame=none,
  aboveskip=0pt,
  belowskip=0pt,
  lineskip=0pt,
}
\lstdefinestyle{promptlistdark}{
  basicstyle=\color{white}\PromptMono,
  columns=fullflexible,
  breaklines=true,
  breakautoindent=false,
  tabsize=2,
  showstringspaces=false,
  frame=none,
  aboveskip=0pt,
  belowskip=0pt,
  lineskip=0pt,
}
\newcommand{\AppendixSnippet}[1]{appendix_snippets/#1}
\newenvironment{PromptSystemBox}[1]{%
  \begin{tcolorbox}[
    width=\linewidth,
    enhanced, arc=2mm, outer arc=2mm,
    colback=black, colframe=black, colupper=white,
    coltitle=white, colbacktitle=black,
    fonttitle=\bfseries\sffamily\footnotesize,
    title={#1}, titlerule=0pt, boxrule=0pt,
    top=1mm, boxsep=3mm,
    before upper={\lstset{linewidth=\linewidth}},
    before skip=\baselineskip, after skip=\baselineskip,
    breakable]
}{\end{tcolorbox}}
\newenvironment{PromptUserBox}[1]{%
  \begin{tcolorbox}[
    width=\linewidth,
    enhanced, arc=2mm, outer arc=2mm,
    colback=white, colframe=black, colupper=black,
    coltitle=black, colbacktitle=white,
    fonttitle=\bfseries\sffamily\small,
    title={#1}, titlerule=0.55pt, titlerule style=black,
    boxrule=0.85pt, top=0.5mm, boxsep=3.5mm, left=3.5mm, right=3.5mm,
    before upper={\lstset{linewidth=\linewidth}},
    before skip=\baselineskip, after skip=\baselineskip,
    breakable]
}{\end{tcolorbox}}
\newenvironment{DefnBox}[1]{%
  \begin{tcolorbox}[
    width=\linewidth,
    enhanced, arc=2mm,
    colback=DefnBoxGray, colframe=black!55,
    fonttitle=\bfseries\footnotesize,
    title={#1},
    boxrule=0.55pt, boxsep=3mm, breakable]
}{\end{tcolorbox}}
\newcommand{\LstPrompt}[1]{\lstinputlisting[style=promptlist]{\AppendixSnippet{#1}}}
\newcommand{\LstPromptDark}[1]{\lstinputlisting[style=promptlistdark]{\AppendixSnippet{#1}}}

\graphicspath{{figures/}{results/figures/}}

\newcommand{\missingfigure}[2]{%
  \fbox{%
    \begin{minipage}[c][0.25\textheight][c]{0.92\linewidth}
      \centering
      Figure asset not found: \texttt{\detokenize{#1}}\\[0.5em]
      #2
    \end{minipage}%
  }%
}
\newcommand{\maybeincludegraphics}[2]{%
  \IfFileExists{#1}{%
    \Description{#2}%
    \includegraphics[width=\linewidth]{#1}%
  }{%
    \IfFileExists{results/#1}{%
      \Description{#2}%
      \includegraphics[width=\linewidth]{results/#1}%
    }{%
      \missingfigure{#1}{#2}%
    }%
  }%
}

\title[Conflict as Structural Stress in LLM Cross-Personality Rewriting]{Conflict as Structural Stress in LLM Cross-Personality Rewriting}

\author{Kaixun Wang}
\affiliation{
  \institution{Southern University of Science and Technology}
  \city{Shenzhen}
  \country{China}
}

\begin{document}

\begin{abstract}
We frame cross-persona conflict as \textbf{structural stress}, not as a uniform ``hybrid vigor'' boost: outcomes \textbf{split by dimension}---automatic novelty (combined) rises with conflict curvature ($\beta_{d^2}{>}0$, Table~\ref{tab:quadratic}), while value and coherence bend downward ($\beta_{d^2}{<}0$), and headline creativity sits between those forces ($\beta_{d^2}{=}{-}0.111$, $p{\approx}4.6{\times}10^{-7}$, $n{=}4620$).
The pre-registered T3-only fit places a quadratic vertex at $d^\ast{\approx}0.47$ on discrete-main T3 rows (Table~\ref{tab:h1}), signalling \textbf{medium} conflict within that subset; pooled binned means of $C_{\mathrm{auto}}$ vs.\ $d_H$ nonetheless trend broadly downward at scale (Figure~\ref{fig:main}), so we do \emph{not} describe headline creativity as a classic rise-then-fall ``inverted-U'' in the aggregate plot.
The pooled sample comprises \textbf{4716} rows (two main arms, mechanism \textbf{60}, multihop 96). Exploratory bins tie rising conflict to higher mean sentiment displacement and non-monotone sentence-order similarity (Figure~\ref{fig:variability}). On the mechanism arm, mean $C_{\mathrm{auto}}$ ranks \textbf{M1}${>}$\textbf{M3}${>}$\textbf{M2}${>}$\textbf{M0} on fifteen high-$d$ cells per mode (descriptive).
Separately, open-source \textbf{Qwen3-4B/8B} replications on the same discrete T3 protocol each show significantly negative quadratic curvature on $C_{\mathrm{auto}}$ vs.\ $d_H$ in the \textbf{same Gaussian mixed-effects specification} as Table~\ref{tab:scale_regression} (genre fixed effects with source-level random intercepts; $n{=}720$ T3 rows per checkpoint: $\beta_{d^2}{=}{-}0.592$ and $-0.505$, respectively, both $p{<}10^{-3}$).
Figure~\ref{fig:scale_inverted_u} is a \textbf{visual supplement only}: it overlays genre-averaged \emph{OLS} predicted curves on the same Qwen rows to sketch shape and hybrid bands; those OLS vertices are not re-used as the primary inferential estimand and need not coincide numerically with the MixedLM $d^\ast$ column in Table~\ref{tab:scale_regression}.
The judge-composite analogue appears in Figure~\ref{fig:scale_inverted_u_judge}.
\end{abstract}

\keywords{AI writing, persona prompting, creativity evaluation, style transfer, structural stress, LLM-as-judge, mixed-effects regression}

\maketitle

\section{Scope of This Report}

Experiments span corpus-scale rewriting with automatic metrics and LLM-as-judge scores, merged for pooled analysis.
Mixed-effects regressions use all rows with finite conflict scalar $d_H$ and finite outcomes on each dependent variable (hence $n{=}4620$ for headline automatic creativity, fewer than the 4716 pooled rows because some rows lack finite $d_H$ or outcomes).
Discrete four-arm summaries (T0--T3 means) use the discrete main arm only, so they are not duplicated alongside the parallel continuous-target arm.

\paragraph{Interpretive frame (structural stress).}
Rather than treating conflict as a homogeneous creativity elevator, we emphasize \textbf{dimensional splitting}: the same stressor can elevate novelty-like signals while eroding value and coherence support for $C_{\mathrm{auto}}{=}N_{\mathrm{auto}}{\cdot}V_{\mathrm{auto}}$, consistent with creativity theories that distinguish generative fluency from evaluative constraint~\cite{boden} and with dual-process accounts that separate rapid associative production from slower monitoring of consistency and judgment~\cite{kahneman}.
Under this reading, an interior optimum near $d^\ast{\approx}0.47$ (Table~\ref{tab:h1}) marks a \textbf{medium-conflict} regime where rewriting remains productive before coherence/value penalties steepen---not an endorsement of maximal persona distance.

\section{Experimental Design}

\noindent\textbf{Appendices.}
Full rewrite and judge prompt templates, automatic-metric formulas, supplementary regression/binning robustness tables, and the exploratory structural-collapse GEE are in Appendices~\ref{app:prompts}--\ref{app:collapse}.

\subsection{Persona Space and Conflict}

Space-H maps each source to $(S,R)\in[-1,1]^2$: increasing $S$ toward $+1$ indicates more analytic, System-2-like prose (toward $-1$, more intuitive and affective); increasing $R$ toward $+1$ indicates a more adventurous register, toward $-1$ a more conservative one. Persona prototypes carry fixed coordinates, and rewrite rows inherit conflict from the source--target pair.

\paragraph{Mapping sources into Space-H.}
The default pipeline uses \textbf{deterministic HuggingFace probes} (no LLM sampling): $S$ clips $(f{-}e)$ where $f$ is a RoBERTa formality probability and $e$ is non-neutral emotion mass from a DistilRoBERTa emotion classifier; $R$ is the dot product of a StyleDistance sentence embedding with a fixed unit axis loaded from precomputed axis vectors (so register contrast is geometric rather than hand-labelled). An alternate configuration uses an LLM coordinate scorer that returns JSON $(S,R)$ on a 1--5 scale, averaged over repeated seeds with spread retained as a stability diagnostic---both backends normalize to $[-1,1]$ before conflict computation.

\paragraph{Reliability and genre structure.}
Coordinates are screened with a \textbf{genre-separation diagnostic}: per-genre means and standard deviations of $S$ and $R$, pairwise Euclidean distances between genre centroids in $(S,R)$, and one-way ANOVA plus Kruskal--Wallis tests for differences in $S$ and $R$ across genres. These checks justify treating Space-H as an ordinal scaffold for conflict rather than as a literal psychometric scale; absolute distances should still be interpreted cautiously.
Figure~\ref{fig:coord_space_h} plots all corpus sources in $(S,R)$ by genre under the same deterministic mapping used for conflict construction.

\begin{figure}[t]
\centering
\maybeincludegraphics{figures/coord_space_h_distribution.png}{Missing figure.}
\caption{Source-level Space-H coordinates $(S,R)$ by genre ($n{=}60$ sources). Axis directions follow the project persona specification: $S$ runs from System-1 (intuitive, affective) at $-1$ to System-2 (analytic, rational) at $+1$; $R$ runs from conservative register at $-1$ to adventurous register at $+1$.}
\label{fig:coord_space_h}
\end{figure}

Scalar conflict $d_H$ and directional $\Delta_H=(\Delta S,\Delta R)$ are computed from these vectors (Euclidean $d_H$ after the project's normalization).

\subsection{Main Arms: Discrete and Continuous Targets}

The discrete main arm uses T0 (identity), T1 (nearest persona), T2 (random persona placebo), and T3 (directed cross-personality).
A parallel continuous-target arm repeats the same structure with continuously sampled T3 targets; both arms enter pooled regressions and genre summaries, while condition means in Table~\ref{tab:condition} use the discrete arm only.
Figure~\ref{fig:d_distribution_discrete_continuous} contrasts the induced T3 distributions of scalar conflict $d_H$: discrete targets concentrate on a small set of corner-persona distances, whereas continuous targets fill the unit disk and yield lower typical $d_H$ (the same comparison yields a handful of descriptive summary statistics in the supplementary materials).

\begin{table}[t]
\centering
\small
\begin{tabular}{lrrrr}
\toprule
Arm & Rows & T0 & T1/T2 & T3 or modes \\
\midrule
Discrete main & 2280 & 120 & 360 / 360 & 1440 T3 \\
Continuous main & 2280 & 120 & 360 / 360 & 1440 T3 \\
Mechanism & 60 & -- & -- & 15 each M0--M3 \\
Multihop & 96 & -- & -- & 96 hop rows \\
\midrule
Total pooled & 4716 & 240 & 1440 / 1440 & 2880 T3, 60 mech, 96 mh.\\
\bottomrule
\end{tabular}
\caption{Row counts by experimental arm. T0--T3 counts are per main arm (two arms). Multihop rows use hop indices rather than T0--T3 labels.}
\label{tab:coverage}
\end{table}

\begin{figure}[t]
\centering
\maybeincludegraphics{figures/d_distribution_discrete_vs_continuous.png}{Missing figure.}
\caption{T3 conflict distance $d_H$ under discrete persona corners versus continuously sampled targets in Space-H ($n{=}1440$ T3 rows per main arm). Overlaid density histograms with a shared bin count.}
\label{fig:d_distribution_discrete_continuous}
\end{figure}

\paragraph{Open-source small-model replication arm.}
Three \textbf{Qwen3} checkpoints (4B, 8B, 14B) are served through a \textbf{vLLM} local deployment and run the \emph{same} discrete main-arm protocol (T0--T3, joint mode, identical prompts and row schema) as the proprietary-generator arms; all downstream metric computation, judging, merging, and analysis use the same pipeline so that cross-scale comparisons isolate generator capacity rather than preprocessing differences.

\section{Measurement Model}

\subsection{Automatic novelty $N_{\mathrm{auto}}$}

$N_{\mathrm{auto}}$ is the equal-weight mean of three \textbf{genre-relative} channels, each mapped to $[0,1]$ after comparing the rewrite to statistics of \emph{source texts in the same genre}: (i)~distinct-2 type-token ratio vs.\ genre mean/variance; (ii)~bigram novelty rate vs.\ typical intra-genre overlap; (iii)~optional SBERT-based $1{-}$cosine distance vs.\ mean pairwise embedding distance within the genre (this channel is informative only when embedding computation is enabled). Z-score--style clipping stabilizes outliers before rebinding to $[0,1]$. The design deliberately removes genre-wide verbosity confounds that would otherwise make poetry appear ``novel'' by construction.

\subsection{Automatic value $V_{\mathrm{auto}}$}

Let entailment be the probability of natural-language inference \emph{entailment} from source premise to generated hypothesis (cross-encoder MNLI). Let coherence be GPT-2 perplexity of the rewrite mapped monotonically to $[0,1]$ (inverse perplexity proxy for fluency). The primary value scalar is the \textbf{arithmetic mean} of entailment and coherence. A geometric mean variant $\sqrt{e\cdot c}$ is computed for robustness tables only.

\subsection{Automatic creativity $C_{\mathrm{auto}}$}

Primary automatic creativity is the \textbf{product}
\[
  C_{\mathrm{auto}} = N_{\mathrm{auto}}\, V_{\mathrm{auto}},
\]
clipped so both factors lie in $[0,1]$ beforehand. Multiplication encodes a conjunctive requirement: high scores need both divergence from the source \emph{and} retention of propositional/fluency support---consistent with definitions that treat creativity as combining originality with acceptable usefulness~\cite{runco}. An additive composite would reward dominance on one axis alone and blur interpretation when novelty and value move against each other under conflict.

Pooling therefore lets novelty and value/coherence move in opposite directions under rising $d_H$, producing the split documented in Table~\ref{tab:quadratic}.

\subsection{LLM-as-judge (validity arm)}

Absolute judging uses a fixed rubric: seven Likert dimensions (novelty, surprise, imagery, fidelity, fluency, consistency, logical completeness) with strict JSON outputs on a 1--5 scale; each generation is scored independently relative to the source and genre. \textbf{Dual judges} are used in this study---DeepSeek V4 Pro and OpenAI GPT-4o---spanning distinct provider families rather than repeated calls to the same endpoint. Per-dimension scores are averaged across judges with successful parses. A pairwise creative-comparison prompt is implemented for head-to-head evaluation with randomized option order, but the merged analysis emphasises the absolute branch for stable row-wise fields. Aggregated judge novelty pools novelty-type dimensions; low correlation with $C_{\mathrm{auto}}$ therefore signals construct divergence between human-like rubric novelty and the automatic NLI$\times$perplexity definition rather than ``failure'' of either side in isolation.

Judge fields are merged into the analysis table for validity checks; they \textbf{do not} replace $C_{\mathrm{auto}}$ as the prespecified primary outcome.

\subsection{Judge validity and construct checks}

Pearson correlations among judge and automatic scores include: judge novelty vs.\ automatic novelty (combined) $r{=}0.528$ (moderate agreement); automatic creativity vs.\ automatic novelty $r{=}0.635$; automatic creativity vs.\ value $r{=}0.668$; \textbf{automatic creativity vs.\ judge novelty} $r{=}0.118$.
The near-zero association between judge-assigned novelty and automatic creativity cautions against treating the two as interchangeable: judge and automatic scores capture \textbf{partially distinct} aspects of rewriting quality, so downstream conclusions should specify which channel they speak to.
Judge fidelity vs.\ automatic novelty is negative ($r{=}{-}0.319$), consistent with a novelty--fidelity tension; judge coherence and fidelity remain positively associated ($r{\approx}0.52$).

To stress-test whether this channel dependence extends beyond linear correlations, we repeated the \textbf{T3-only} scale-style regressions used alongside Table~\ref{tab:scale_regression}, replacing $C_{\mathrm{auto}}$ with the judge-based composite $C_{\mathrm{judge}}$ (computed from merged judge scores in the same analysis pipeline).
Each fit uses $C_{\cdot} \sim d_H + d_H^2 + C(\mathrm{genre})$ with a source-level random intercept on T3 rows (Qwen checkpoints: $n{=}720$ each; proprietary OpenAI runs pooled across two model assignments: $n{=}1440$ T3 discrete rows in one fit, as in Table~\ref{tab:judge_auto_divergence}).

\begin{table}[t]
\centering
\small
\begin{tabular}{lrrrrrr}
\toprule
Generator & $\beta_{d^2}^{\mathrm{auto}}$ & $p(d^2)$ & Sig. & $\beta_{d^2}^{\mathrm{judge}}$ & $p(d^2)$ & Sig. \\
\midrule
Qwen3-4B  & $-0.592$ & $2.9{\times}10^{-5}$ & Yes & $-0.417$ & $0.188$ & No \\
Qwen3-8B  & $-0.505$ & $2.2{\times}10^{-4}$ & Yes & $+0.242$ & $0.500$ & No \\
Qwen3-14B & $-0.050$ & $0.719$               & No  & $+0.194$ & $0.604$ & No \\
GPT-4o    & $-0.618$ & $7.0{\times}10^{-10}$ & Yes & $+1.007$ & $1.8{\times}10^{-5}$ & Yes \\
\bottomrule
\end{tabular}
\caption{T3-only quadratic curvature for automatic vs.\ judge creativity by generator (mixed-effects with genre fixed effects and source random intercepts). GPT-4o shows significant \emph{negative} curvature on $C_{\mathrm{auto}}$ and significant \emph{positive} curvature on $C_{\mathrm{judge}}$ simultaneously---the clearest within-arm signature that the two scoring channels are not measuring the same construct.}
\label{tab:judge_auto_divergence}
\end{table}

On the Qwen3 arms---where the same downstream judging stack is shared across checkpoints---$C_{\mathrm{judge}}$ need not track $C_{\mathrm{auto}}$.
For \textbf{Qwen3-4B}, the $d_H^2$ slopes share direction (both negative) but only the automatic channel is conventionally significant; the judge channel is not ($p{\approx}0.19$).
For \textbf{Qwen3-8B}, the signs \textbf{invert}: $C_{\mathrm{auto}}$ keeps the familiar negative curvature, whereas $C_{\mathrm{judge}}$ turns positive, so at this model's capacity frontier high-conflict rewrites can remain attractive to human-aligned rubric scores even as the NLI-and-perplexity channel deteriorates.
Neither channel should be read as a substitute for the other: the split is direct evidence of \textbf{construct difference}, not ``noise'' in a single gold standard.
For \textbf{Qwen3-14B}, neither quadratic term is significant on either channel in this slice, consistent with weaker separation at scale.

The proprietary \textbf{GPT-4o} arm is the starkest case (Table~\ref{tab:judge_auto_divergence}): $\beta_{d^2}^{\mathrm{auto}}{=}{-}0.618$ ($p{\approx}7.0{\times}10^{-10}$) and $\beta_{d^2}^{\mathrm{judge}}{=}{+}1.007$ ($p{\approx}1.8{\times}10^{-5}$).
Both are highly significant yet point in opposite directions, which is difficult to reconcile with a single latent ``creativity'' unless one allows the instruments to weight different stress responses.
We interpret this as systematic \textbf{construct divergence}: the automatic composite primarily penalizes \emph{semantic drift} away from the source under rising conflict, whereas the judge composite rewards \emph{expressive surprise} encoded in the rubric---both are nominally about ``creative'' rewrites, but they encode different definitions of what counts as a gain.
This nuance also helps explain the very low headline alignment between judge novelty and $C_{\mathrm{auto}}$ ($r{=}0.118$; Figure~\ref{fig:aux}, left panel): the two channels need not move together when stress pulls lexical--entailment fidelity and human-judged novelty in opposite directions.

\paragraph{Surprise dimension along $d_H$.}
The rubric \emph{surprise} score (Likert mean$/5$ on the same merged judge block) can be regressed with the \textbf{identical} T3-only mixed specification as $C_{\mathrm{auto}}$ and $C_{\mathrm{judge}}$: $\mathrm{surprise}{\sim} d_H{+}d_H^2{+}C(\mathrm{genre})$ with a source random intercept on the \textbf{same} $n{=}1440$ proprietary T3 discrete rows as Table~\ref{tab:judge_auto_divergence} (half GPT-4o, half GPT-4o mini).
Table~\ref{tab:surprise_curvature} shows that \textbf{surprise} carries a large positive quadratic coefficient ($\beta_{d^2}{\approx}{+}1.89$, $p{<}10^{-8}$), closely tracking $C_{\mathrm{judge}}$ ($\beta_{d^2}{\approx}{+}1.01$) and opposite $C_{\mathrm{auto}}$ ($\beta_{d^2}{\approx}{-}0.62$).
Further splitting the rubric shows that the surprise dimension's positive curvature is \textbf{almost twice} that of the headline judge composite and an order of magnitude tighter in $p$-value, so the high-conflict \emph{lift} on $C_{\mathrm{judge}}$ concentrates overwhelmingly on human-rated \emph{surprise} rather than on a uniform shift across all seven Likert legs.
In short: high conflict raises judged quality chiefly by activating perceived surprise, which sharpens the construct-divergence story with a concrete rubric source.

\begin{table}[t]
\centering
\small
\begin{tabular}{lrr}
\toprule
Outcome (T3 pooled, $n{=}1440$) & $\beta_{d^2}$ & $p(d^2)$ \\
\midrule
Surprise$/5$ (rubric mean) & $+1.891$ & $8.1{\times}10^{-9}$ \\
$C_{\mathrm{auto}}$ & $-0.618$ & $7.0{\times}10^{-10}$ \\
$C_{\mathrm{judge}}$ & $+1.007$ & $1.8{\times}10^{-5}$ \\
\bottomrule
\end{tabular}
\caption{T3-only quadratic curvature for rubric surprise vs.\ headline composites (Gaussian mixed-effects with genre fixed effects and source random intercepts).}
\label{tab:surprise_curvature}
\end{table}

\begin{table}[t]
\centering
\small
\begin{tabular}{lrrrrr}
\toprule
Condition & Creativity & Novelty & Value & NLI & Coherence \\
\midrule
T0 & 0.2317 & 0.3334 & 0.7044 & 0.7486 & 0.6602 \\
T1 & 0.3197 & 0.4176 & 0.7677 & 0.7759 & 0.7596 \\
T2 & 0.3120 & 0.4231 & 0.7408 & 0.7153 & 0.7662 \\
T3 & 0.3124 & 0.4232 & 0.7409 & 0.7144 & 0.7675 \\
\midrule
\multicolumn{6}{@{}l@{}}{\footnotesize\emph{Open-source Qwen3 generators, T3 condition only (same discrete protocol as T0--T3 block above).}} \\
Qwen3-4B & 0.2449 & 0.3695 & 0.6735 & 0.5106 & 0.8364 \\
Qwen3-8B & 0.2543 & 0.3870 & 0.6664 & 0.5190 & 0.8138 \\
Qwen3-14B & 0.2508 & 0.3968 & 0.6429 & 0.5096 & 0.7763 \\
\midrule
M0 & 0.1615 & 0.4415 & 0.3879 & 0.4147 & 0.3611 \\
M1 & 0.2044 & 0.4471 & 0.4771 & 0.5652 & 0.3890 \\
M2 & 0.1908 & 0.4583 & 0.4305 & 0.5372 & 0.3238 \\
M3 & 0.2018 & 0.4013 & 0.5084 & 0.5305 & 0.4863 \\
\bottomrule
\end{tabular}
\caption{Condition means. T0--T3: discrete main arm only. M0--M3: mechanism arm, fifteen independent generations per mode on top-$d$ cells. For the open-source Qwen3 rows, lower T3 creativity versus proprietary T3 tracks mainly \textbf{value} and \textbf{NLI} shortfalls rather than novelty, consistent with weaker semantic fidelity in smaller generators.}
\label{tab:condition}
\end{table}

\section{Descriptive Results}

\subsection{Genre-Level Patterns}

Pooling discrete and continuous main arms (4560 rows) yields the following means.

\begin{table}[t]
\centering
\small
\begin{tabular}{lrrrrr}
\toprule
Genre & Creativity & Novelty & Value & NLI & Coherence \\
\midrule
Academic & 0.2838 & 0.3913 & 0.7342 & 0.6879 & 0.7805 \\
Essay & 0.3434 & 0.4466 & 0.7730 & 0.8078 & 0.7383 \\
Narrative & 0.2593 & 0.3628 & 0.7174 & 0.5547 & 0.8801 \\
Poetry & 0.3385 & 0.4520 & 0.7519 & 0.9054 & 0.5984 \\
\bottomrule
\end{tabular}
\caption{Genre means pooling discrete and continuous main arms (4560 rows).}
\label{tab:genre}
\end{table}

\begin{figure*}[t]
\centering
\begin{minipage}[t]{0.49\textwidth}
\centering
\maybeincludegraphics{figures/inverted_u_creativity_auto.png}{Missing figure.}
\end{minipage}\hfill
\begin{minipage}[t]{0.49\textwidth}
\centering
\maybeincludegraphics{figures/directional_field_creativity_auto.png}{Missing figure.}
\end{minipage}
\caption{Binned mean $C_{\mathrm{auto}}$ vs.\ $d_H$ and the directional field in $(\Delta S,\Delta R)$. Headline creativity trends downward with $d_H$ in pooled binned means (left); regressions still estimate negative $d_H^2$ on $C_{\mathrm{auto}}$ over the mixed sample alongside positive novelty curvature (Table~\ref{tab:quadratic}). Right: interactive $\Delta S$--$\Delta R$ structure.}
\label{fig:main}
\end{figure*}

\begin{figure*}[t]
\centering
\begin{minipage}[t]{0.49\textwidth}
\centering
\maybeincludegraphics{figures/judge_auto_validity_matrix.png}{Missing figure.}
\end{minipage}\hfill
\begin{minipage}[t]{0.49\textwidth}
\centering
\maybeincludegraphics{figures/t3_d_H_distribution_main.png}{Missing figure.}
\end{minipage}
\caption{Left: Pearson correlations among judge and automatic scores. Right: $d_H$ distribution on discrete-main T3 rows, contextualizing the H1 vertex relative to observed conflict mass.}
\label{fig:aux}
\end{figure*}

\begin{figure}[t]
\centering
\maybeincludegraphics{figures/generation_variability_sentiment_kendall.png}{Missing figure.}
\caption{Exploratory \textbf{generation variability} vs.\ conflict: eight equal-width $d_H$ bins over the pooled sample. \emph{Top:} mean absolute sentiment displacement ($[0,2]$). \emph{Bottom:} mean sentence-order preservation score ($[0,1]$, higher $\approx$ closer to source order).}
\label{fig:variability}
\end{figure}

\section{Conflict Intensity}

\subsection{Regression specification}

The pooled quadratic models regress each outcome on linear and quadratic conflict terms plus categorical genre and experimental condition:
\[
  y \sim d_H + d_H^2 + \sum_g \mathbf{1}_{\mathrm{genre}=g} + \sum_c \mathbf{1}_{\mathrm{condition}=c}.
\]
Observations are rewrite-level rows; \textbf{random intercepts group by source text}, reflecting repeated generations from the same corpus item without nesting random effects in persona identity. Condition absorbs arm contrasts (T0--T3, mechanism modes, multihop structure as coded); persona distance enters through $d_H$ and directional regressors rather than as an extra random classification. Mixed linear models are used when variance-component estimation is numerically stable; otherwise the same fixed-effects structure is estimated by OLS. In sparse designs the estimated source-level variance may collapse toward zero, so $\beta_{d^2}$ should be read as a pooled association that partly marginalizes within-source correlation rather than as a nested hierarchical causal estimand.
Directional regressions swap $(d_H,d_H^2)$ for $(\Delta S,\Delta R)$ plus squares, product, and Euclidean $|\Delta_H|$, retaining the same genre/condition fixed effects and source-level intercept grouping.

\subsection{Pooled Quadratic Results}

Mixed-effects fits surface the dimensional split: \textbf{negative} $d_H^2$ on headline creativity, value, and coherence versus \textbf{positive} $d_H^2$ on novelty (combined), over $n{=}4620$ rows with valid $d_H$ and finite outcomes on each metric (Table~\ref{tab:quadratic}).
Together with the T3-only vertex $d^\ast{\approx}0.47$ (Table~\ref{tab:h1}), one can read a medium-conflict regime as comparatively favourable \emph{within that subset}, in line with classic arousal--performance tradeoffs~\cite{yerkes}; Figure~\ref{fig:main} nonetheless shows pooled binned $C_{\mathrm{auto}}$ declining with $d_H$, so the vertex should not be equated with a visible global inverted-U in the aggregate bins.

\paragraph{Scale dependence of the inverted-U.}
Table~\ref{tab:scale_regression} fixes the same quadratic specification on automatic creativity ($C_{\mathrm{auto}}\sim d_H+d_H^2+C(\mathrm{genre})$) with \textbf{source-level random intercepts}, fit \textbf{separately per metrics file} on \textbf{T3 discrete-target rows only}.
Each open-source Qwen checkpoint file contributes $n{=}720$ T3 rows (one generator model per file).
The proprietary slice pools \textbf{all} T3 discrete-target rows for the two OpenAI generator assignments in a \emph{single} mixed-effects fit ($n{=}1440$: half GPT-4o, half GPT-4o mini), matching Table~\ref{tab:judge_auto_divergence}. This is not ``half the corpus'' and not a discrete-vs-continuous split: continuous-target T3 rows for these generators are not included in that proprietary discrete pool.
Optimisers use L-BFGS-B by default; Powell is used when the Hessian is singular (typical on the narrow Qwen subsets).

\begin{table}[t]
\centering
\small
\begin{tabular}{lrrrl}
\toprule
Generator (T3) & $\beta_{d^2}$ & $p(d^2)$ & $d^\ast$ & Sig.\ ($p{<}0.05$) \\
\midrule
Qwen3-4B & $-0.5925$ & $2.89{\times}10^{-5}$ & 0.4811 & Yes \\
Qwen3-8B & $-0.5052$ & $2.24{\times}10^{-4}$ & 0.5319 & Yes \\
Qwen3-14B & $-0.0501$ & 0.7185 & 0.0185 & No \\
GPT-4o (pooled) & $-0.6181$ & $6.95{\times}10^{-10}$ & 0.4954 & Yes \\
\bottomrule
\end{tabular}
\caption{Quadratic curvature for $C_{\mathrm{auto}}$ vs.\ $d_H$ by generator (T3-only discrete rows; mixed-effects with source intercepts). Qwen checkpoints: $n{=}720$ each. \textbf{GPT-4o (pooled)}: $n{=}1440$ (half GPT-4o, half GPT-4o mini) in one fit, matching Table~\ref{tab:judge_auto_divergence}.}
\label{tab:scale_regression}
\end{table}

\begin{table}[t]
\centering
\small
\begin{tabular}{lrrr}
\toprule
Metric & $n$ & $\beta_{d^2}$ & $p(d^2)$ \\
\midrule
Creativity auto & 4620 & $-0.1111$ & $4.58{\times}10^{-7}$ \\
Creativity geom & 4620 & $-0.1479$ & $2.76{\times}10^{-8}$ \\
Value auto & 4620 & $-0.3901$ & $7.22{\times}10^{-17}$ \\
Coherence auto & 4620 & $-0.3266$ & $5.37{\times}10^{-11}$ \\
Novelty (combined) & 4620 & $0.0908$ & $1.90{\times}10^{-6}$ \\
\bottomrule
\end{tabular}
\caption{Quadratic terms from pooled mixed-effects regressions on headline outcomes and components. Singular random-effect covariance warnings can occur at numerical boundaries.}
\label{tab:quadratic}
\end{table}

\subsection{Pre-Registered H1 Check}

The fitted vertex $d^\ast{\approx}0.47$ operationalizes \textbf{medium} conflict within discrete-main T3 mass (empirical q03--q97 shown); it lies outside the narrower pre-registered vertex band yet inside the sample-mass diagnostic, so claims should reference both the band logic and the descriptive optimum. This interior optimum is estimated on T3 rows only and coexists with monotonic-looking pooled binned means vs.\ $d_H$ (Figure~\ref{fig:main}).

\begin{table}[t]
\centering
\small
\begin{tabular}{lr}
\toprule
H1 diagnostic & Value \\
\midrule
T3 main rows & 1440 \\
$\beta_d$ & 0.3265 \\
$\beta_{d^2}$ & $-0.3481$ \\
$p(d^2)$ & $6.51{\times}10^{-4}$ \\
Fitted vertex $d$ & 0.4689 \\
Empirical $d$ q03 / q97 & 0.2948 / 0.7971 \\
Vertex inside pre-reg band & False \\
Vertex inside sample mass & True \\
Pooled main $p(d^2)$ & $4.58{\times}10^{-7}$ \\
\bottomrule
\end{tabular}
\caption{Pre-registered H1 diagnostic for automatic creativity.}
\label{tab:h1}
\end{table}

The open-source replications sharpen where ``medium conflict'' sits in $d_H$: \textbf{Qwen3-4B} concentrates useful mass roughly in $d_H\in[0.30,0.63]$, while \textbf{Qwen3-8B} widens slightly to $d_H\in[0.32,0.74]$ around the same quadratic peaks (Figure~\ref{fig:scale_inverted_u}), offering generator-specific hybrid operating bands beyond the pooled GPT-only vertex in Table~\ref{tab:h1}. The same construction applied to the judge composite $C_{\mathrm{judge}}$ is shown in Figure~\ref{fig:scale_inverted_u_judge} for direct visual comparison with the automatic channel.

\begin{figure}[t]
\centering
\maybeincludegraphics{figures/scale_inverted_u.png}{Missing figure.}
\caption{T3-only genre-averaged OLS predictions of $C_{\mathrm{auto}}$ vs.\ $d_H$ for Qwen3 checkpoints (4B/8B/14B). Vertical dotted lines mark fitted vertices for 4B/8B; shaded bands sketch approximate hybrid intervals (4B: $[0.30,0.63]$; 8B: $[0.32,0.74]$).}
\label{fig:scale_inverted_u}
\end{figure}

\begin{figure}[t]
\centering
\maybeincludegraphics{figures/scale_inverted_u_judge.png}{Missing figure.}
\caption{Same axes and curve construction as Figure~\ref{fig:scale_inverted_u}, but with mean predicted $C_{\mathrm{judge}}$ (genre-averaged OLS) on the vertical axis. Contrasts the rubric-based composite against the automatic $C_{\mathrm{auto}}$ trajectories on identical Qwen3 T3 rows.}
\label{fig:scale_inverted_u_judge}
\end{figure}

\subsection{Conflict direction}

\begin{table}[t]
\centering
\small
\begin{tabular}{lrr}
\toprule
Term & Coefficient & $p$ \\
\midrule
$\Delta S$ & 0.0164 & $3.57{\times}10^{-43}$ \\
$\Delta R$ & 0.0028 & 0.0127 \\
$\Delta S^2$ & $-0.0018$ & 0.607 \\
$\Delta R^2$ & $-0.0143$ & $3.67{\times}10^{-3}$ \\
$\Delta S\Delta R$ & 0.0075 & $1.93{\times}10^{-11}$ \\
$|\Delta|$ & 0.0190 & 0.0343 \\
\bottomrule
\end{tabular}
\caption{Directional regression for automatic creativity.}
\label{tab:direction}
\end{table}

The same directional structure reappears when we restrict to the \textbf{open-source Qwen3-4B/8B} T3 slices: the interaction term $\Delta S{\times}\Delta R$ remains the dominant nonlinear channel (highly significant on 8B, jointly informative on 4B alongside negative $\Delta R^2$ curvature), matching the pooled GPT-family pattern in Table~\ref{tab:direction} and supporting a shared interpretation that oblique moves in $(\Delta S,\Delta R)$---not pure radial distance---shape automatic creativity.

\subsection{Causal contrasts}

Bucketed contrasts pool \textbf{both} main metric files (hence larger per-bucket $n$ than a discrete-only table).

Across all three conflict buckets, mean $C_{\mathrm{auto}}$ at T3 lies below both T1 and T2 (Table~\ref{tab:causal}), and both T3--T1 and T3--T2 contrasts are negative or near zero. At this pooled scale, \textbf{directed cross-persona targeting (T3) does not outperform} nearest-neighbour (T1) or random persona assignment (T2) on automatic creativity; if anything, it trails both---a pattern that should temper expectations that ``designed'' cross-persona moves automatically beat placebo-style persona swaps.

\begin{table}[t]
\centering
\small
\begin{tabular}{lrrrrrrrr}
\toprule
Bucket & $n_{T1}$ & $n_{T2}$ & $n_{T3}$ & T3--T1 & T3--T2 & mean T1 & mean T2 & mean T3 \\
\midrule
Low & 228 & 68 & 1004 & $-0.0023$ & $-0.0009$ & 0.3127 & 0.3113 & 0.3104 \\
Mid & 468 & 240 & 841 & $-0.0165$ & $-0.0061$ & 0.3211 & 0.3107 & 0.3046 \\
High & 24 & 412 & 1035 & $-0.0523$ & $-0.0078$ & 0.3591 & 0.3146 & 0.3068 \\
\bottomrule
\end{tabular}
\caption{Bucketed causal contrasts on automatic creativity.}
\label{tab:causal}
\end{table}

On the open-source \textbf{Qwen3} replication metrics, the pre-registered \textbf{H5} (quadratic curvature in JSD-based divergence vs.\ $d_H$) and \textbf{H6} (T3 vs.\ T2 contrasts within binned conflict) are again \emph{non-significant}, paralleling the proprietary pattern above: directed T3 persona choice does not deliver incremental gains over the T2 placebo within fixed $d_H$ bins.
This is a substantive empirical finding rather than a methodological defect, so we report it here alongside the causal buckets rather than under limitations.

\paragraph{Structural tail check (appendix).}
The automatic-channel deterioration under high conflict is already conveyed by $C_{\mathrm{auto}}$, NLI, and coherence slopes in the main tables.
For readers who want a \textbf{binary} articulation of the same stress axis, Appendix~\ref{app:collapse} fits a clustered binomial GEE on a joint low-NLI-and-low-coherence indicator (sample-relative $q{=}0.2$ cut-offs on the same $n{=}1440$ proprietary T3 pool); marginal fitted risk rises steeply at high $d_H$.
We keep that material in an appendix because it introduces GEE machinery and sample-defined thresholds that sit awkwardly with the main mixed-effects narrative, and because the qualitative message largely overlaps what is already visible in the automatic scores.

\section{Mechanism and Multihop Branches}

\subsection{Mechanism Experiment}

Mechanism generations restrict to the fifteen highest-$d$ discrete-main cells and cross four informational modes (M0--M3); each mode has fifteen rows (Table~\ref{tab:condition}).

\paragraph{Which instructional mechanism helps most?}
Pooling those sixty rows, mean automatic creativity orders \textbf{M1} (direct rewrite, highest mean $C_{\mathrm{auto}}$) ${>}$ \textbf{M3} (constraint-forward rewrite) ${>}$ \textbf{M2} (two-stage) ${>}$ \textbf{M0} (control). This addresses the project-book theme that different conflict-\emph{instruction} bundles yield different creative profiles: unconstrained cross-persona rewriting (M1) best preserves $C_{\mathrm{auto}}$ here, tightly constrained rewriting (M3) is second, explicit decomposition (M2) sits mid-pack, and the passive control (M0) is lowest---all on the same high-$d$ cells.
With $n{=}15$ per mode this ordering is \textbf{descriptive}; formal mechanism-only inference would warrant dedicated models with pair-level blocking.

Table~\ref{tab:mechanism} restricts to rows with finite $d_H$ and outcome, takes the upper third of $d_H$ \emph{globally}, then aggregates by mode. The M0--M3 lines describe mechanism modes inside that high-$d$ slice; the joint-mode line summarizes continuous-main rows under that label in the same slice ($n{=}1470$), not the sixty-row mechanism-only sample.

\begin{table}[t]
\centering
\small
\begin{tabular}{lrrr}
\toprule
Mode & Mean creativity & Std. & Count \\
\midrule
M0 & 0.1615 & 0.0701 & 15 \\
M1 & 0.2044 & 0.0793 & 15 \\
M2 & 0.1908 & 0.0859 & 15 \\
M3 & 0.2018 & 0.0674 & 15 \\
Joint & 0.3099 & 0.0765 & 1470 \\
\bottomrule
\end{tabular}
\caption{Mechanism-related summary in the global high-$d$ slice (see text).}
\label{tab:mechanism}
\end{table}

\subsection{Multihop Prediction Check}

\begin{table}[t]
\centering
\small
\begin{tabular}{llrrr}
\toprule
Genre & Bin & Multihop & Direct T3 & Delta \\
\midrule
Academic & Low & 0.3244 & 0.2941 & 0.0303 \\
Academic & Mid & 0.2134 & 0.2830 & $-0.0696$ \\
Academic & High & 0.2553 & 0.2736 & $-0.0182$ \\
Essay & Low & 0.3968 & 0.3412 & 0.0557 \\
Essay & Mid & 0.3948 & 0.3527 & 0.0421 \\
Essay & High & 0.4142 & 0.3429 & 0.0713 \\
Narrative & Low & 0.2695 & 0.2706 & $-0.0011$ \\
Narrative & Mid & 0.2781 & 0.2515 & 0.0266 \\
Narrative & High & 0.2733 & 0.2619 & 0.0115 \\
Poetry & Low & 0.3333 & 0.3298 & 0.0035 \\
Poetry & Mid & 0.4294 & 0.3404 & 0.0890 \\
Poetry & High & 0.3704 & 0.3512 & 0.0192 \\
\bottomrule
\end{tabular}
\caption{Multihop vs.\ direct T3 mean automatic creativity by genre and $d$ bin.}
\label{tab:multihop}
\end{table}

Multihop effects are \textbf{genre-dependent}. Essay rows show multihop strictly above direct T3 in every $d$ bin (all positive deltas), suggesting smoother persona paths help that register here. Academic rows reverse under mid conflict (delta ${\approx}{-}0.07$), while low/high bins are mixed or modest---so gradual hops are not universally preferable and should be evaluated per genre and conflict band.

\section{Limitations}

\begin{itemize}\itemsep2pt
\item Pooled regressions mix discrete and continuous main arms with mechanism and multihop rows; interpret coefficients as pooled associations, not as a single pre-registered estimand.
\item Mixed-effects fits may emit singular-covariance or convergence warnings.
\item Causal buckets and mechanism summaries depend on binning and filtering choices in the analysis implementation.
\item Mechanism arm contrasts use fifteen draws per mode on fixed high-$d$ cells; generalization to other persona pairs requires replication.
\item Sentiment/Kendall variability curves (Figure~\ref{fig:variability}) are exploratory and share the pooled binning choice (eight equal-width $d_H$ bins).
\item For smaller checkpoints (especially Qwen3), \textbf{genre-averaged OLS} curves in Figure~\ref{fig:scale_inverted_u} are illustrative; their nominal standard errors do not incorporate the same hierarchical structure as the MixedLM rows in Table~\ref{tab:scale_regression} and can be \textbf{anti-conservative} when read as inferential objects rather than as sketches of shape.
\item The ``\textbf{capacity threshold}'' language for Qwen3-8B (judge curvature turning positive while automatic curvature stays negative) is a \textbf{post hoc descriptive} reading of two fitted slopes at different model sizes, not an independently identified causal capacity parameter.
\item Appendix~\ref{app:collapse} reports a binomial GEE on sample-defined collapse cut-offs (not a calibrated GLMM); treat as exploratory triangulation, not as an absolute failure probability.
\end{itemize}

\section{Conclusion}

Cross-persona conflict behaves as \textbf{structural stress}: it does \emph{not} uniformly lift creativity; instead it reallocates mass across automatic subscores---novelty curvature turns positive while value and coherence turn sharply negative (Table~\ref{tab:quadratic}). The discrete-main T3-only quadratic places a vertex near $d^\ast{\approx}0.47$ (Table~\ref{tab:h1}), while pooled binned means of $C_{\mathrm{auto}}$ vs.\ $d_H$ trend downward at scale (Figure~\ref{fig:main}); together these caution against equating ``medium conflict'' with a visibly inverted-U aggregate trajectory for headline creativity.
Exploratory variability traces (Figure~\ref{fig:variability}) show sentiment displacement rising into mid-to-high conflict bins and the sentence-order score dipping through mid bins, consistent with more aggressive reordering co-occurring with the value/coherence erosion channel.
Directionally, $\Delta S$, $\Delta R^2$, and $\Delta S\Delta R$ remain central (Table~\ref{tab:direction}). On mechanism prompts over shared high-$d$ cells, mean $C_{\mathrm{auto}}$ orders M1${>}$M3${>}$M2${>}$M0 (Table~\ref{tab:condition}), a descriptive answer to ``which bundle works best'' within this arm.
Open-source Qwen3-4B/8B runs reproduce the same qualitative story at smaller scale: a significant inverted-U for $C_{\mathrm{auto}}$ vs.\ $d_H$ with optima near $d^\ast{\approx}0.46$--$0.53$ and practical hybrid bands $d_H\in[0.30,0.63]$ (4B) / $[0.32,0.74]$ (8B) (Table~\ref{tab:scale_regression}; Figure~\ref{fig:scale_inverted_u}; judge-channel analogue Figure~\ref{fig:scale_inverted_u_judge}).
Treat all claims as conditional on branch (main vs.\ mechanism vs.\ multihop), sampling mode, and the narrow mechanism sample size.

For practitioners, the choice of evaluation metric matters as much as the choice of conflict distance. On the automatic channel ($C_{\mathrm{auto}}$, which penalizes semantic drift), medium conflict near $d^\ast{\approx}0.47$ is comparatively favourable and high conflict degrades quality---making it the right metric for fidelity-sensitive tasks such as academic editing or translation. On the judge channel ($C_{\mathrm{judge}}$, which rewards human-perceived surprise), the curvature reverses: GPT-4o shows $\beta_{d^2}^{\mathrm{judge}}{=}{+}1.007$ ($p{<}10^{-4}$), meaning higher conflict produces more positively-evaluated rewrites by human-like rubrics---making it the right metric for expressive tasks such as creative writing or advertising copy. The two channels are not interchangeable; downstream design should select the metric that matches the intended use case before tuning conflict distance.

\appendix
% Appendix: boxed ``system / user'' prompt cards (cf. ranking-prompt appendix style).
\begingroup\sloppy

\section{Appendix A: Prompt templates (boxed layout)}
\label{app:prompts}

\noindent\textbf{Implementation note.}
Full discrete persona bodies, coordinate-scorer definitions, and paraphrase variants live in the project's configuration; generation, mechanism, judging, and experiment orchestration follow the accompanying open-source implementation in the repository.

\subsection{A.1 Main arms: T0--T3 rewrites}

\subsubsection{T0 --- Identity (measurement floor)}
\begin{DefnBox}{Pipeline note}
No LLM call: the stored row is the source text verbatim with $d_H{=}0$. No system/user messages.
\end{DefnBox}

\subsubsection{T1--T3 discrete --- Joint generator (\texttt{rewrite\_joint})}
\begin{PromptSystemBox}{Rewrite / Joint --- System input}
\footnotesize
\texttt{<persona.system\_prompt>} from the persona configuration (primary or paraphrase variant), plus the length line \texttt{Target length: L--H words.}
\end{PromptSystemBox}
\begin{PromptUserBox}{Rewrite / Joint --- User input}
\LstPrompt{rewrite_user_template.txt}
\end{PromptUserBox}

\subsubsection{T3 continuous target --- Synthetic persona header}
Continuous sampling builds a lightweight synthetic persona header; numeric $S,R$ are filled per draw. The user message is the same joint wrapper as above.
\begin{PromptSystemBox}{Rewrite / Continuous T3 --- System input (pattern)}
\bgroup\color{white}\footnotesize\ttfamily\raggedright
Adopt this continuous style target in Space-H: S=+0.000, R=+0.000.\par
Lean toward rational and adventurous expression proportionally to magnitudes.\par
\egroup
\end{PromptSystemBox}
\begin{PromptUserBox}{Rewrite / Continuous T3 --- User input}
\LstPrompt{rewrite_user_template.txt}
\end{PromptUserBox}

\subsection{A.2 Mechanism arms: M0--M3}

\subsubsection{M1 --- Joint (same as main joint)}
One-shot \texttt{rewrite\_joint}: system $=$ persona $+$ length; user $=$ shared wrapper (Section~A.1).

\subsubsection{M2 --- Sequential paraphrase then persona}
\begin{PromptSystemBox}{Mechanism M2 / Stage 1 --- System input}
\bgroup\color{white}\footnotesize\ttfamily\raggedright
You are a faithful paraphraser. You restate the source's propositional content in neutral, unadorned English prose. Do not add, remove, or alter any commitment of the source. Keep the length within 90--110\% of the source.\par
\egroup
\end{PromptSystemBox}
\begin{PromptUserBox}{Mechanism M2 / Stage 1 --- User input}
\LstPrompt{paraphrase_user.txt}
\end{PromptUserBox}
\begin{PromptSystemBox}{Mechanism M2 / Stage 2 --- System input}
\footnotesize
\texttt{<persona.system\_prompt>} $+$ length constraint (same as joint).
\end{PromptSystemBox}
\begin{PromptUserBox}{Mechanism M2 / Stage 2 --- User input}
\LstPrompt{m2_stage2_user_template.txt}
\end{PromptUserBox}

\subsubsection{M3 --- Checklist-constrained joint}
\begin{PromptSystemBox}{Mechanism M3 / Checklist extraction --- System input}
\bgroup\color{white}\footnotesize\ttfamily\raggedright
You are a careful reading-comprehension assistant. From the passage below, extract a numbered checklist of its propositional commitments (factual claims, causal links, and concrete imagery) that any faithful rewrite must preserve. Be concise.\par
\egroup
\end{PromptSystemBox}
\begin{PromptUserBox}{Mechanism M3 / Checklist extraction --- User input}
\LstPrompt{checklist_user.txt}
\end{PromptUserBox}
\begin{PromptSystemBox}{Mechanism M3 / Persona + checklist --- System input}
\footnotesize
\texttt{<persona.system\_prompt>} $+$ length, followed by the injected checklist banner (trimmed to a word budget).
\end{PromptSystemBox}
\begin{PromptUserBox}{Mechanism M3 --- Checklist banner appended to system block}
\LstPrompt{m3_inject_header.txt}
\end{PromptUserBox}
\begin{PromptUserBox}{Mechanism M3 --- User message (same as joint)}
\LstPrompt{rewrite_user_template.txt}
\end{PromptUserBox}

\subsubsection{M0 --- Checklist control (token-pressure placebo)}
Checklist is still extracted for auditing, but the text appended to the persona system block is a length-matched placeholder list (control path).
\begin{PromptUserBox}{Mechanism M0 --- Control banner (system block append)}
\LstPrompt{m0_control_header.txt}
\end{PromptUserBox}
\begin{PromptUserBox}{Mechanism M0 --- User message (same as joint)}
\LstPrompt{rewrite_user_template.txt}
\end{PromptUserBox}

\subsection{A.3 Judge prompts (absolute and pairwise)}

\subsubsection{Absolute rubric (seven Likert dimensions)}
\begin{PromptSystemBox}{Judge / Absolute --- System input}
\texttt{You are a strict JSON-only literary judge.}
\end{PromptSystemBox}
\begin{PromptUserBox}{Judge / Absolute --- User input}
\LstPrompt{judge_absolute_user.txt}
\end{PromptUserBox}
\begin{DefnBox}{JSON parsing and aggregation}
Judge completions are parsed as JSON. Each Likert dimension is clipped to $[1,5]$; per-dimension scores are averaged across judges with valid JSON. A generous token ceiling is used for JSON-shaped outputs.
\end{DefnBox}

\subsubsection{Pairwise creative comparison}
\begin{PromptSystemBox}{Judge / Pairwise --- System input}
\texttt{You are a strict JSON-only literary judge.}
\end{PromptSystemBox}
\begin{PromptUserBox}{Judge / Pairwise --- User input}
\LstPrompt{judge_pairwise_user.txt}
\end{PromptUserBox}
\begin{DefnBox}{Pairwise scoring note}
\texttt{REWRITE\_A}/\texttt{B} presentation order is randomised; signed scores are mapped back to ``A vs.\ B'' after accounting for order.
\end{DefnBox}

\subsection{A.4 Coordinate scorer (LLM backend only)}
When the coordinate backend is LLM-based, sources are rated with the scorer prompt in the persona configuration (integer $S,R\in[-5,5]$, JSON with reasons, then normalise to $[-1,1]$). The runs analysed in the main text use the deterministic HF backend; the LLM configuration branch still documents axis semantics for reviewers.

\section{Appendix B: Automatic metric definitions}
\label{app:metrics}

\begin{DefnBox}{B.1 Novelty N-auto (combined channel)}
Let $c^{\mathrm{d2}}_{\mathrm{rel}}$, $c^{\mathrm{bg}}_{\mathrm{rel}}$, $c^{\mathrm{emb}}_{\mathrm{rel}}$ be genre-relative scores in $[0,1]$. Then
\[
  N_{\mathrm{auto}} = \tfrac{1}{3}\bigl(c^{\mathrm{d2}}_{\mathrm{rel}} + c^{\mathrm{bg}}_{\mathrm{rel}} + c^{\mathrm{emb}}_{\mathrm{rel}}\bigr).
\]
Distinct-2: $z$-score vs.\ genre source mean/std, clip $z\in[-3,3]$, rebound $(z/3+1)/2$. Bigram novelty: offset vs.\ genre baseline with heuristic spread $\sigma{=}0.1$, same clip/rebound. Embedding channel: neutral $0.5$ when SBERT is off; else normalised vs.\ intra-genre mean pairwise distance with clipping.
\end{DefnBox}

\begin{DefnBox}{B.2 Value $V_{\mathrm{auto}}$}
Coherence maps GPT-2 perplexity through a monotone pivot map (default pivot 80). Let $e$ be NLI entailment in $[0,1]$ and $c$ coherence in $[0,1]$. Primary value is the arithmetic mean $(e+c)/2$. Geometric value uses $\sqrt{e\cdot c}$ with a zero guard.
\end{DefnBox}

\begin{DefnBox}{B.3 Creativity variants}
$C_{\mathrm{auto}}=N_{\mathrm{auto}}\cdot V_{\mathrm{arith}}$ (clipped factors); $C_{\mathrm{auto,geom}}=N_{\mathrm{auto}}\cdot V_{\mathrm{geom}}$. Winsorised diagnostics multiply winsorised novelty by arithmetic value in the same analysis protocol as the main tables.
\end{DefnBox}

\section{Appendix C: Statistical robustness}
\label{app:robust}

\subsection{C.1 Pooled quadratic: OLS vs.\ mixed-effects}
Rows match Table~\ref{tab:quadratic}: discrete main, continuous main, and mechanism arms only ($n{=}4620$ with finite $d_H$ and outcomes). Mixed models add a random intercept by source text; optimiser order matches the main analysis script defaults.

\begin{table}[t]
\centering
\begin{tcolorbox}[
  width=\linewidth,
  enhanced, arc=2mm, colback=white, colframe=black,
  title={\texorpdfstring{\textbf{Table C.1.} OLS vs.\ MixedLM on $\beta_{d^2}$ (conflict quadratic)}{\textbf{Table C.1.} OLS vs. MixedLM on beta d2 (conflict quadratic)}},
  fonttitle=\bfseries\small, titlerule=0.55pt, boxrule=0.95pt, boxsep=3mm]
\small
\noindent\resizebox{\linewidth}{!}{%
\begin{tabular}{@{}l l r r r@{}}
\toprule
Metric & Model & $n$ & $\beta_{d^2}$ & $p(d^2)$ \\
\midrule
Creativity auto & MixedLM (L-BFGS, REML) & 4620 & $-0.1111$ & $4.58{\times}10^{-7}$ \\
Creativity auto & OLS (same FE)        & 4620 & $-0.1072$ & $5.97{\times}10^{-5}$ \\
Novelty (combined) & MixedLM (L-BFGS, REML) & 4620 & $0.0908$ & $1.90{\times}10^{-6}$ \\
Novelty (combined) & OLS (same FE)        & 4620 & $0.1143$ & $1.37{\times}10^{-5}$ \\
\bottomrule
\end{tabular}%
}
\end{tcolorbox}
\caption{Quadratic conflict coefficient: sign and significance are stable across OLS and mixed-effects estimators on the pooled main$+$mechanism sample.}
\label{tab:appendix-robust-reg}
\end{table}

\subsection{C.2 T3--T2 gaps under alternative $d_H$ binning}
\begin{table}[t]
\centering
\begin{tcolorbox}[
  width=\linewidth,
  enhanced, arc=2mm, colback=white, colframe=black,
  title={\texorpdfstring{\textbf{Table C.2.} Mean $C_{\mathrm{auto}}$ difference (T3 minus T2) by bucket}{\textbf{Table C.2.} Mean C auto difference (T3 minus T2) by bucket}},
  fonttitle=\bfseries\small, titlerule=0.55pt, boxrule=0.95pt, boxsep=3mm]
\small
\noindent\resizebox{\linewidth}{!}{%
\begin{tabular}{@{}l r r r@{}}
\toprule
Binning scheme & low: T3$-$T2 & mid: T3$-$T2 & high: T3$-$T2 \\
\midrule
Quantile tertiles & $-0.0009$ & $-0.0061$ & $-0.0078$ \\
Equal-width 3 bins & ---       & $-0.0102$ & $-0.0017$ \\
\bottomrule
\end{tabular}%
}
\end{tcolorbox}
\caption{Mean $C_{\mathrm{auto}}$ gaps (T3 minus T2) inside conflict buckets on the same $n{=}4620$ pool; negative values match the main text's ``T3 does not beat T2'' pattern under both binning rules where defined.}
\label{tab:appendix-robust-bin}
\end{table}

\section{Appendix D: Structural collapse risk (exploratory GEE)}
\label{app:collapse}

This appendix reproduces the exploratory \textbf{structural collapse} analysis deferred from the main text's conflict-intensity section.

\subsection{D.1 Indicator and GEE specification}
\begin{DefnBox}{D.1 Collapse event (sample-relative)}
$\mathrm{collapse}{=}1$ iff $\mathrm{NLI}{<}\tau_{\mathrm{NLI}}$ and $\mathrm{coherence}_{\mathrm{auto}}{<}\tau_{\mathrm{coh}}$, with $\tau_{\cdot}$ at the \textbf{same-sample 20th percentiles} on the $n{=}1440$ proprietary T3 discrete rows (same pool as judge-divergence and surprise regressions).
\end{DefnBox}
\noindent
This yields $\tau_{\mathrm{NLI}}{\approx}0.412$, $\tau_{\mathrm{coh}}{\approx}0.628$, and baseline $\hat P(\mathrm{collapse}){\approx}1.9\%$.
Because statsmodels does not expose a standard logit GLMM with Gaussian random intercepts alongside our Gaussian mixed models, we fit a \textbf{binomial GEE} with clustering by source text and exchangeable working correlation (fallback independence); the estimand is marginal (population-averaged) risk.
The genre-inclusive specification $\mathrm{collapse}{\sim}d{+}d^2{+}C(\mathrm{genre})$ shows quasi-separation in narrative/poetry strata; we therefore emphasise the \emph{no-genre} robustness fit for marginal probabilities.
On a fine $d_H$ grid, implied $P(\mathrm{collapse})$ rises from $\approx 3.5\%$ at $d_H{\approx}0.22$ to $\approx 88\%$ at $d_H{\approx}0.88$, with the steepest increase near $d_H{\approx}0.83$; the quadratic term remains significant at $p{\approx}0.003$.
Table~\ref{tab:appendix_collapse_gee} summarises the headline marginal diagnostics.
Full log-odds coefficient tables, the genre-inclusive sensitivity fit, Hessian/convergence metadata, and a median-threshold ($q{=}0.5$) variant are archived with the supplementary regression output for this appendix.

\begin{table}[t]
\centering
\begin{tcolorbox}[
  width=\linewidth,
  enhanced, arc=2mm, colback=white, colframe=black,
  title={\texorpdfstring{\textbf{Table D.1.} Marginal $P(\mathrm{collapse})$ diagnostics (no-$C(\mathrm{genre})$ GEE, $q{=}0.2$)}{\textbf{Table D.1.} Marginal P collapse diagnostics}},
  fonttitle=\bfseries\small, titlerule=0.55pt, boxrule=0.95pt, boxsep=3mm]
\small
\noindent\resizebox{\linewidth}{!}{%
\begin{tabular}{@{}l r@{}}
\toprule
Diagnostic & Value \\
\midrule
Rows ($n$) & 1440 \\
Baseline $\hat P(\mathrm{collapse})$ & ${\approx}1.9\%$ \\
$\hat P(\mathrm{collapse})$ at $d_H{\approx}0.22$ & ${\approx}3.5\%$ \\
$\hat P(\mathrm{collapse})$ at $d_H{\approx}0.88$ & ${\approx}88\%$ \\
$d_H$ at maximum slope of $\hat P$ & ${\approx}0.83$ \\
$p(d_H^2)$ (quadratic term) & ${\approx}0.0032$ \\
\bottomrule
\end{tabular}%
}
\end{tcolorbox}
\caption{Structural-collapse risk: marginal fitted probabilities on $d_H$ (binomial GEE with clustering by source). Log-odds coefficients are omitted because they become large whenever $\hat P$ spans nearly the unit interval; full coefficient tables accompany the supplementary materials.}
\label{tab:appendix_collapse_gee}
\end{table}

\endgroup


\begin{thebibliography}{9}
\bibitem{kahneman}
D.~Kahneman.
\newblock \textit{Thinking, Fast and Slow}.
\newblock Farrar, Straus and Giroux, 2011.

\bibitem{boden}
M.~Boden.
\newblock \textit{The Creative Mind: Myths and Mechanisms}.
\newblock Routledge, 2004.

\bibitem{runco}
M.~Runco and G.~Jaeger.
\newblock The standard definition of creativity.
\newblock \textit{Creativity Research Journal} 24, 1 (2012), 92--96.

\bibitem{yerkes}
R.~Yerkes and J.~Dodson.
\newblock The relation of strength of stimulus to rapidity of habit-formation.
\newblock \textit{Journal of Comparative Neurology and Psychology} 18 (1908), 459--482.
\end{thebibliography}

\end{document}
